\chapter{Results \& Performance Analysis}
\label{chap:results}

\section{Fall Detection Calibration \& Kinematic Validation}
To evaluate the classification fidelity of the MPU6050 fall detection algorithm and calibrate optimal detection thresholds, empirical trials were conducted comparing traumatic fall scenarios against common Activities of Daily Living (ADL). 

Table~\ref{tab:fall_calibration} presents the recorded peak tri-axial accelerations, total vector magnitude ($A_{\text{total}}$), angular velocities ($G_x, G_y, G_z$), and Euler orientation angles ($\text{Pitch}, \text{Roll}, \text{Yaw}$) across controlled physical trials.

\begin{table}[htbp]
    \centering
    \caption{Kinematic Sensor Calibration Dataset for Fall Events vs. Daily Living Activities}
    \label{tab:fall_calibration}
    \begin{tabularx}{\textwidth}{l c c c c c c c c c c}
        \toprule
        \textbf{Activity / Trial} & \textbf{$A_x$ ($g$)} & \textbf{$A_y$ ($g$)} & \textbf{$A_z$ ($g$)} & \textbf{$A_{\text{total}}$ ($g$)} & \textbf{$G_x$} & \textbf{$G_y$} & \textbf{$G_z$} & \textbf{Pitch} & \textbf{Roll} & \textbf{Yaw} \\
        \midrule
        \textbf{Forward Fall} & $2.45$ & $0.85$ & $1.90$ & $\mathbf{3.21}$ & $4.12$ & $3.85$ & $1.20$ & $1.42$ & $0.31$ & $0.55$ \\
        \textbf{Backward Fall} & $1.20$ & $0.65$ & $2.85$ & $\mathbf{3.16}$ & $3.95$ & $4.50$ & $0.95$ & $1.51$ & $0.22$ & $0.48$ \\
        \textbf{Lateral Fall (Left)} & $0.95$ & $2.60$ & $1.50$ & $\mathbf{3.15}$ & $2.80$ & $1.10$ & $4.20$ & $0.35$ & $1.48$ & $0.62$ \\
        \textbf{Lateral Fall (Right)} & $0.88$ & $2.75$ & $1.42$ & $\mathbf{3.21}$ & $2.95$ & $1.25$ & $4.35$ & $0.32$ & $1.52$ & $0.59$ \\
        \textbf{Vertical Collapse} & $0.62$ & $0.55$ & $2.95$ & $\mathbf{3.06}$ & $1.85$ & $2.10$ & $1.15$ & $1.38$ & $0.40$ & $0.30$ \\
        \midrule
        \textbf{Prayer Prostration} & $0.60$ & $0.20$ & $0.40$ & $\mathbf{0.75}$ & $0.51$ & $0.56$ & $0.50$ & $0.61$ & $0.24$ & $0.35$ \\
        \textbf{Rapid Sitting} & $0.35$ & $0.25$ & $1.65$ & $\mathbf{1.70}$ & $1.10$ & $0.85$ & $0.40$ & $0.25$ & $0.15$ & $0.20$ \\
        \textbf{Brisk Walking} & $0.45$ & $0.30$ & $1.25$ & $\mathbf{1.36}$ & $0.95$ & $1.20$ & $0.80$ & $0.18$ & $0.12$ & $0.45$ \\
        \textbf{Stair Descent} & $0.55$ & $0.40$ & $1.80$ & $\mathbf{1.92}$ & $1.35$ & $1.10$ & $0.65$ & $0.22$ & $0.18$ & $0.30$ \\
        \textbf{Bending to Floor} & $0.70$ & $0.30$ & $0.65$ & $\mathbf{1.00}$ & $0.80$ & $0.75$ & $0.45$ & $0.85$ & $0.20$ & $0.15$ \\
        \bottomrule
    \end{tabularx}
\end{table}

As shown in Table~\ref{tab:fall_calibration}, all true traumatic fall impacts generated $A_{\text{total}} \ge 3.06g$, significantly surpassing the calibrated threshold of $\text{Th}_{\text{impact}} = 2.80g$. In contrast, aggressive daily activities such as rapid sitting and stair descent peaked at $1.70g$ and $1.92g$ respectively. Controlled prayer prostration demonstrated slow kinematic rotation without the preceding $<0.5g$ freefall signature ($A_{\text{total}} = 0.75g$), confirming that the multi-stage algorithm effectively eliminates false triggers during routine non-emergency physical movement.

In addition, across 25 intentional dropped-device trials, the 10-second acoustic/visual cancellation grace period achieved a $100\%$ false-alarm suppression rate when cancelled by the operator within the allocated window.

\section{Wireless Telemetry \& Failover Latency Analysis}
The communication reliability of the primary LoRaWAN (AU915, SF7, $125\,\text{kHz}$) link and secondary Wi-Fi fallback was characterized under simulated industrial noise. 

\begin{itemize}
    \item \textbf{LoRaWAN Airtime \& Packet Delivery:} The 4-byte binary telemetry payload achieved a total physical airtime of $T_{\text{packet}} = 36.1\,\text{ms}$. Across 500 transmission trials distributed across the facility, the primary LoRa link achieved a packet delivery success rate of $98.4\%$, with an average uplink latency of $142\,\text{ms} \pm 18\,\text{ms}$.
    \item \textbf{Autonomous Wi-Fi Failover:} Under forced LoRaWAN concentrator disconnection (simulating gateway failure or deep structural blockage), the client node executed 2 unacknowledged retries ($2.1\,\text{s}$ timeout) before successfully establishing a Wi-Fi socket connection and dispatching the emergency packet. The total end-to-end failover dispatch latency was measured at $3.96\,\text{s} \pm 0.35\,\text{s}$, well within the critical 10-second emergency response SLA.
\end{itemize}

\section{BLE Indoor Positioning \& Gateway Attenuation Profiling}

\subsection{Gateway RSSI Attenuation Characterization}
\label{sec:gateway_rssi_comparison}
Empirical signal attenuation profiles were gathered across the facility to benchmark the RF sensitivity and stability of the Minew G1, Mokosmart MKGW3, and Mokosmart LW003-B gateways. Table~\ref{tab:gateway_rssi_comparison} presents the measured average RSSI ($\text{dBm}$) at stepped distances from $0\,\text{m}$ to $7\,\text{m}$.

\begin{table}[htbp]
    \centering
    \caption{Empirical Gateway RSSI Attenuation Across Stepped Distances}
    \label{tab:gateway_rssi_comparison}
    \begin{tabularx}{0.85\textwidth}{c c c c}
        \toprule
        \textbf{Distance ($d$, meters)} & \textbf{Minew G1 ($\text{dBm}$)} & \textbf{Mokosmart MKGW3 ($\text{dBm}$)} & \textbf{Mokosmart LW003-B ($\text{dBm}$)} \\
        \midrule
        $\mathbf{0\,\text{m}}$ (Near Field) & $-30\,\text{dBm}$ & $-60\,\text{dBm}$ & $-58\,\text{dBm}$ \\
        $\mathbf{1\,\text{m}}$ & $-50\,\text{dBm}$ & $-75\,\text{dBm}$ & $-72\,\text{dBm}$ \\
        $\mathbf{2\,\text{m}}$ & $-65\,\text{dBm}$ & $-80\,\text{dBm}$ & $-78\,\text{dBm}$ \\
        $\mathbf{3\,\text{m}}$ & $-70\,\text{dBm}$ & $-85\,\text{dBm}$ & $-82\,\text{dBm}$ \\
        $\mathbf{4\,\text{m}}$ & $-70\,\text{dBm}$ & $-85\,\text{dBm}$ & $-84\,\text{dBm}$ \\
        $\mathbf{5\,\text{m}}$ & $-65\,\text{dBm}$ & $-80\,\text{dBm}$ & $-80\,\text{dBm}$ \\
        $\mathbf{6\,\text{m}}$ & $-55\,\text{dBm}$ & $-75\,\text{dBm}$ & $-76\,\text{dBm}$ \\
        $\mathbf{7\,\text{m}}$ & $-50\,\text{dBm}$ & $-75\,\text{dBm}$ & $-74\,\text{dBm}$ \\
        \bottomrule
    \end{tabularx}
\end{table}

As shown in Table~\ref{tab:gateway_rssi_comparison}, the Minew G1 exhibited higher receiver sensitivity ($-50\,\text{dBm}$ at $1\,\text{m}$ vs. $-75\,\text{dBm}$ for Mokosmart). However, at distances beyond $4\,\text{m}$, constructive multipath reflections from metallic structural walls caused the RSSI to artificially rise (reaching $-50\,\text{dBm}$ at $7\,\text{m}$). This non-monotonic behavior demonstrates the fundamental limitation of raw geometric trilateration in enclosed industrial environments.

\subsection{Localization Algorithm Accuracy Comparison}
To resolve multipath anomalies, four localization methodologies were implemented and benchmarked across a 4-gateway layout (\texttt{ble\_positioning\_4gw.py}). Table~\ref{tab:positioning_performance} summarizes the resulting positioning performance.

\begin{table}[htbp]
    \centering
    \caption{Indoor Positioning Performance Benchmark Across Localization Methods}
    \label{tab:positioning_performance}
    \begin{tabularx}{\textwidth}{l c c X}
        \toprule
        \textbf{Localization Algorithm} & \textbf{Mean Coordinate Error} & \textbf{Room Classification Accuracy} & \textbf{Computational Complexity / Stability} \\
        \midrule
        \textbf{Raw RSSI Trilateration} & $3.85\,\text{m} \pm 1.42\,\text{m}$ & $62.4\%$ & High coordinate jitter; ill-conditioned in multipath. \\
        \textbf{Kalman-Filtered Trilateration} & $2.15\,\text{m} \pm 0.68\,\text{m}$ & $78.6\%$ & Smooth trajectory; reduces transient noise spikes. \\
        \textbf{Weighted Centroid (WCL)} & $2.40\,\text{m} \pm 0.85\,\text{m}$ & $81.2\%$ & Stable, bounded inside convex hull of gateways. \\
        \textbf{Fuzzy Logic Room Classifier} & $\mathbf{N/A \text{ (Symbolic)}}$ & $\mathbf{94.8\%}$ & Highly robust; eliminates coordinate jumping. \\
        \bottomrule
    \end{tabularx}
\end{table}

The experimental data confirms that while discrete Kalman filtering substantially reduces raw coordinate error (from $3.85\,\text{m}$ to $2.15\,\text{m}$), the Fuzzy Inference System mapping to predefined Room Zones achieved superior real-world reliability ($94.8\%$ room classification accuracy), ensuring clear location tags in emergency dispatches.

\section{Power Budget Kinetics \& Server UPS Reliability}

\subsection{Client Battery Discharge Characterization}
Continuous discharge experiments were conducted on the wearable client comparing $600\,\text{mAh}$ vs. $1100\,\text{mAh}$ LiPo cells under active $100\,\text{ms}$ BLE advertising and periodic LoRaWAN keepalive transmissions:

\begin{itemize}
    \item \textbf{600 mAh LiPo Cell:} Initiated at 9:30 AM ($4.20\,\text{V}$), the device depleted to hardware cutoff ($3.40\,\text{V}$) at 3:15 PM, yielding a total operational runtime of $5\,\text{hours } 45\,\text{minutes}$. Linear fuel gauge interpolation reported a premature shutdown percentage of $35\%$.
    \item \textbf{1100 mAh LiPo Cell:} Operating under identical workloads, the $1100\,\text{mAh}$ cell achieved $>24.5\,\text{hours}$ of continuous operation, providing an industrial Safety Factor ($\text{SF}$) of:
    \begin{equation}
    \text{SF} = \frac{\text{Operational Runtime}}{\text{Standard Shift Duration}} = \frac{24.5\,\text{h}}{12.0\,\text{h}} = \mathbf{2.04}
    \end{equation}
    \item \textbf{Non-Linear Voltage Remapping:} Reworking the battery estimation algorithm to follow the true electro-chemical discharge curve (mapping $4.20\,\text{V} \rightarrow 100\%$, $3.70\,\text{V} \rightarrow 50\%$, and $3.30\,\text{V} \rightarrow 0\%$) eliminated premature cutoff reporting, providing monotonic, high-accuracy battery telemetry.
\end{itemize}

\subsection{Server UPS Mains Failure Endurance}
The central Raspberry Pi 5 server connected to the dual-cell Waveshare UPS HAT was evaluated under simulated total mains power loss:
\begin{itemize}
    \item \textbf{Mains Disconnection Timestamp:} 4:10 PM (Day 1).
    \item \textbf{Automated Low-Battery SMS Warning:} Triggered at 1:34 AM (Day 2) when cell voltage reached $3.30\,\text{V}$.
    \item \textbf{Safe Automated Shutdown:} Executed cleanly by \texttt{ups-shutdown.service} at 2:02 AM (Day 2) at $3.20\,\text{V}$ ($0.0\%$ capacity).
    \item \textbf{Total Server Backup Endurance:} $\mathbf{9\,\text{hours } 52\,\text{minutes}}$ of continuous multithreaded server, MQTT broker, and web GUI operation during building power blackout.
\end{itemize}

\section{Hardware Diagnostic Analysis: Raspberry Pi 5 RP1 PMIC Lockup}
During testing of early edge server assemblies, a critical non-recoverable boot failure occurred on the Raspberry Pi 5 platform. Connecting a UART-to-USB console cable ($115200\,\text{baud}$) to the dedicated debug port yielded the boot log:
\begin{lstlisting}[language=bash]
0.97 RPi: BOOTSYS release VERSION:2226a853 DATE: 2025/12/08
0.98 AON_RESET: 00000003 PM_RSTS 00001000
0.99 POWER_OFF_ON_HALT: 0 WAIT_FOR_POWER_BUTTON 0 power-on-reset 1
0.99 RP1_BOOT chip ID: 0x0000cf00
1.00 I2C error @ 800058a2
\end{lstlisting}

\textbf{Root Cause Analysis:} The expected RP1 chip identifier is \texttt{0x20001927}. The incorrect ID (\texttt{0x0000cf00}) combined with the I2C bus error indicated that the external UPS HAT was back-powering and pulling the shared I2C telemetry lines high \textit{before} the Pi 5's internal PMIC ($3.3\,\text{V}$ rail) had completed its power-on sequence. This created bus contention that locked the RP1 I/O controller. The issue was permanently mitigated by integrating Schottky isolation diodes on the I2C lines and introducing a $200\,\text{ms}$ power-sequencing delay circuit.

\section{Bill of Materials (BOM) \& Unit Financial Analysis}
Table~\ref{tab:procurement_bom} outlines the complete Bill of Materials for the fully functional prototype platform.

\begin{table}[htbp]
    \centering
    \caption{Emergency Response System Prototype Procurement \& Bill of Materials}
    \label{tab:procurement_bom}
    \begin{tabularx}{\textwidth}{l c c r l}
        \toprule
        \textbf{Component Description} & \textbf{Qty} & \textbf{Unit Cost (AUD)} & \textbf{Total (AUD)} & \textbf{Supplier} \\
        \midrule
        Raspberry Pi Pico 2WH & 3 & \$14.30 & \$42.90 & Core Electronics \\
        Waveshare LoRaWAN HAT for Pico (SX1262) & 4 & \$24.15 & \$96.60 & Core Electronics \\
        Seeed Studio XIAO nRF52840 BLE Module & 4 & \$19.25 & \$77.00 & Core Electronics \\
        InvenSense MPU6050 6-DOF IMU & 3 & \$4.55 & \$13.65 & Core Electronics \\
        LiPo Battery 1S 3.7V 1100mAh & 2 & \$8.65 & \$17.30 & Core Electronics \\
        Waveshare Pico UPS Module (Type B) & 3 & \$21.90 & \$65.70 & Core Electronics \\
        915MHz LoRa Antenna \& RF Cable & 2 & \$20.50 & \$41.00 & Core Electronics \\
        Emergency Push Button \& Tactile Switches & 4 & \$1.30 & \$5.20 & Core Electronics \\
        Raspberry Pi 5 Model B (4GB RAM) & 1 & \$179.55 & \$179.55 & Core Electronics \\
        Waveshare SX1302/SX1303 LoRa Gateway HAT & 1 & \$159.95 & \$159.95 & Core Electronics \\
        Waveshare Pi 5 UPS HAT \& 18650 Cells & 1 & \$99.40 & \$99.40 & Core Electronics \\
        Minew G1 BLE-to-Wi-Fi Gateway & 1 & \$150.00 & \$150.00 & DigiKey Australia \\
        Veroboard, Connectors, Wiring \& Hardware & 1 & \$17.85 & \$17.85 & Core Electronics \\
        \midrule
        \textbf{Total Prototype Hardware Expenditure} & & & \textbf{\$966.10} & \textit{Excl. Shipping} \\
        \bottomrule
    \end{tabularx}
\end{table}

At scale ($>500\,\text{units}$), migrating to a single consolidated 4-layer PCB with SMT components is projected to reduce the per-wearable hardware BOM from $\$87.50\,\text{AUD}$ to under $\$32.00\,\text{AUD}$, establishing a highly competitive price-to-performance ratio.

